\chapter*{Úvod}
\phantomsection
\addcontentsline{toc}{chapter}{Úvod}

Cílem práce je analýza možností úniku akustických vibrací (například, lidská řeč) prostřednictvím stávajících optických datových infrastruktur. 
Teoretická část práce bude obsahovat rozbor možností realizace optických vláknových mikrofonů a srovnání jejich základních parametrů. První kapitola bude obsahovat přehled optického vlákna, princip jeho funkce a možnosti jeho využití pro přenos médií. 
Další kapitola se věnuje optickým vláknovým mikrofonům, principu funkce a jejich konstrukcím. V~rámci semestrální práce bude provedena analýza vhodnosti použití optických vláknových mikrofonů pro implementace odposlechu. Rovněž bude proveden návrh testovacího pracoviště včetně měření s komerčním optickým vláknovým mikrofonem.

Závěrečná práce bude zaměřena na realizaci vlastního optického vláknového mikrofonu umožňujícího měření na stávajících datových rozvodech, bude provedeno měření, analýza a vyhodnocení získaných dat.

Detekce a analýza akustických vln je pole výzkumu s velkým počtem aplikací, které v posledních několika desetiletích hodně rostly. Význam akustické vlnové technologie jak v moderní vědě, tak v aplikačním inženýrství rychle roste. Aplikace se rozšířily do různých oblasti, jako jsou navigace (sledování tlaku a rychlosti, měření hloubky/profilování mořského dna), nedestruktivní vyhodnocení struktur pomocí detekce emitovaných ultrazvuků (detekce a lokalizace trhlin, vnitřní napětí, mikro vysídlení a další), detekce fázových přechodů ve vědeckém výzkumu, ve zdravotnictví (lékařské zobrazování a diagnóza), může se používat pro vojenské a privátní záležitostí (jako odposlech). Senzory s optickými vlákny odstraní všechna omezení akustického snímání, které mají klasické mikrofony a přináší další výhody. Mezi jejich výhody patří: malá velikost, nízká cena, vysoká citlivost a schopnost měřit širokou škálu různých vstupních parametrů. Systémy s optickými vlákny mají velkou šířku pásma, nejsou omezeny na velikost a tvar, jsou robustní a schopné multiplexingu. Díky těmto vlastnostem jsou optická vlákna velmi dobrá a perspektivní technologie pro akustické senzory. První systémy byly popsány v roce 1977 Bucaro ét al. a Cole a kol. \cite{21, 23}. Jejich předběžná práce označila začátek rozvoje technologie optického akustického snímání. V současné době v této oblasti lze nalézt mnoho děl, ze kterých vychází tyto následující interferometrické konfigurace: Machová-Zehnderová \cite{24}, Michelsonová \cite{25}, Sagnacová \cite{26}, a Fabryová-Perotová \cite{27}, zatímco ostatní prozkoumávají možnosti nabízené FBG senzory \cite{28}.

